\chapter{Injection beam loss monitoring system}\label{sec:Injection bea loss monitoring system}

\section{Required performance}
SuperKEKB加速器においてビームは10usごとにリングを一周する。ターンバイターンのロスを測定するには高い時間分解能が要求される。またコリメータ付近に設置されるセンサーから読み出しボードまでは数百mのケーブルで繋がっており、振動の減衰が起きる。そのため取得信号のゲインが十分に大きいことも必要である。コリメータ付近は放射線レベルが高いので、放射線耐性を持つことも求められる。


CLAWSの信号は1Hzに規格化された入射キック信号と同期して取得される。取得された信号は数百メートルのケーブルを経由して地上の電源棟に送られる。電源棟にはレシーバー基板があり、信号を整える役割を持つ。そのあとに信号はオシロスコープによって読み出される。オシロスコープは読み出し用のPCにより制御されており、トリガー発砲から23ms秒間の波形データが取得される。

\section{Data correction, processing, recording}
23msの生波形データは約100万行のデータとなり、保存するには向かない量である。そこで以下の量を再構成してデータを記録する。
各ターンごとのビームロス量。生波形に対して閾値を設定し、各ターンごとに積分を実行。
各入射ごとの総ビームロス量。各ターンごとのビームロス量の和。
各入射後とのビームロス量の1msまでの和と、1ms以降の和。
各入射ごとのロス継続時間の重心。
以上の量をpvとしてSuperKEKBのデータ保存システムであるArchiver上に記録・保存する。これにより誰にでもアクセスが可能なデータとなる。